\documentclass[12pt,a4paper]{article}

\usepackage[utf8]{inputenc}
\usepackage[T1]{fontenc}
\usepackage{graphicx}
\usepackage{amsmath}
\usepackage{hyperref}
\usepackage{booktabs}
\usepackage{xcolor}
\usepackage{geometry}
\usepackage{float}
\geometry{margin=2.5cm}
\usepackage[skip=0.8em]{parskip}

\title{3D Space Partitioning Viewer\\[0.5em]
       \large Computer Graphics (E016712A)}
\author{Antonio Romero Mart\'in \and Siebe Ternest}
\date{\today}

\begin{document}

\maketitle
\newpage
\tableofcontents
\newpage

% ---------------------------------------------------------------------------
\section{Introduction}
% ---------------------------------------------------------------------------

Efficiently organizing millions of 3D points is a fundamental challenge in
computer graphics and spatial computing. Space partitioning data structures
address this by hierarchically subdividing 3D space, enabling fast queries such
as ray intersection, frustum culling, and nearest-neighbour search. As LiDAR
scanners become ubiquitous in autonomous driving, urban mapping, and robotics,
the choice of partitioning structure has direct consequences for rendering
performance, memory usage, and query latency at scales ranging from thousands
to tens of millions of points.

This report documents the design decisions, implementation challenges, and
comparative analysis of a 3D space partitioning viewer built for LiDAR point
clouds. Three hierarchical data structures were implemented and compared:
an \textbf{Octree}, a \textbf{k-d Tree}, and a \textbf{Binary Space
Partitioning (BSP) Tree}. All three were evaluated on three real-world LiDAR
datasets -- a small indoor corridor scan, a mid-size UFO-shaped structure, and
the city of Hasselt ($\sim$19.6 million points) -- covering a wide range of point
densities and spatial configurations.

The viewer is built with \textbf{Three.js} (WebGL), which provides scene-graph
management, instanced geometry, and camera controls out of the box, allowing
the implementation effort to focus on algorithm design and rendering
architecture rather than low-level shader infrastructure.
The project is structured as a single-page web application served by Vite,
with all three algorithms in a shared \texttt{src/algorithms/} module and all
visualization logic in \texttt{src/viewer/}.

Section~\ref{sec:input} describes the input data pipeline and PCD file parsing.
Section~\ref{sec:algorithms} presents each algorithm's design decisions,
complexity, and memory characteristics. Section~\ref{sec:viewer} documents the
viewer and rendering architecture. Section~\ref{sec:analysis} provides a
comparative analysis across build cost, query performance, and suitability for
different scene types.

% ---------------------------------------------------------------------------
\section{Reading the Input Data}
\label{sec:input}
\textit{(Both authors)}
% ---------------------------------------------------------------------------

Input datasets are PCD (Point Cloud Data) files from LiDAR scanners.
The module \texttt{src/loaders/pcdLoader.js} implements a \textbf{fast-path
parser} for the layout shared by all three bundled datasets:
\texttt{FIELDS~x~y~z}, \texttt{SIZE~4~4~4}, \texttt{TYPE~F~F~F},
\texttt{DATA~binary}. The parser reads the ASCII header up to the
\texttt{DATA~binary\textbackslash n} boundary marker and then wraps the
remaining raw bytes directly in a \texttt{Float32Array} view, bypassing the
field-by-field \texttt{DataView.getFloat32()} reads that Three.js
\texttt{PCDLoader} would otherwise perform. For Hasselt ($\sim$19.6 million
points, 235\,MB) this eliminates $\sim$58.8 million individual
\texttt{DataView} calls, reducing the post-download parsing phase from several
seconds to milliseconds. ASCII files and binary PCD with extra fields or
non-float types fall back to Three.js \texttt{PCDLoader} automatically.

The module unifies three loading paths behind a single interface:

\begin{itemize}
  \item \textbf{Bundled datasets} (UFO, Corridor, Hasselt) -- served as static
        assets by the development server; the URL is fetched via XHR (native
        progress events) and the resulting \texttt{ArrayBuffer} is handed to
        the fast-path parser.
  \item \textbf{Arbitrary URL} -- the control panel exposes a URL input field;
        any \texttt{.pcd} URL entered by the user follows the same XHR path and
        is cached by URL string after the first load.
  \item \textbf{Local files} -- \texttt{.pcd} files dropped onto the viewer or
        selected via the file picker arrive as browser \texttt{File} objects;
        \texttt{File.arrayBuffer()} is used to obtain the raw bytes directly,
        which are then passed to the same parser. Local files are cached under
        the key \texttt{\_\_local\_\_:<filename>}, so the user can switch to a
        different dataset and return to the local file instantly without
        re-reading it.
\end{itemize}

Once a point cloud is loaded, \texttt{datasetService} performs two passes over
the position buffer. First, a \textbf{Z-range scan} finds the global minimum and
maximum Z values using a direct typed-array loop rather than
\texttt{Box3.setFromObject}, avoiding Three.js method-call overhead for each of
the 19.6 million points. Second, a \textbf{colorization pass} maps each point's
Z coordinate linearly to an HSL hue in the range
$[0.7\text{ (blue)},\, 0.0\text{ (red)}]$ and writes it into a
\texttt{Float32Array} color attribute. This per-point color is used as the
baseline appearance for all three algorithms. All loaded datasets are \textbf{cached after the first load}: the colorized
\texttt{Points} object, the bounding-box extents, and the camera target are
stored -- URL-based datasets under their URL string, local files under
\texttt{\_\_local\_\_:<filename>} -- so switching between algorithms, adjusting
depth, or returning to a previously loaded dataset never triggers a reload or a
second colorization pass over millions of points.

\subsection{Datasets}

Three binary PCD files were used throughout this project as the primary
evaluation datasets; the viewer also accepts any user-supplied \texttt{.pcd}
file via drag-and-drop or URL input.
Table~\ref{tab:datasets} summarises the key properties of the three bundled
datasets, derived directly from the file headers and bounding-box statistics.

\begin{table}[H]
\centering
\begin{tabular}{@{}lrrrrrr@{}}
\toprule
\textbf{Dataset} & \textbf{Points} & \textbf{Size} &
  \textbf{X span (m)} & \textbf{Y span (m)} & \textbf{Z span (m)} \\
\midrule
Corridor & 1{,}622{,}607  & 19.5\,MB  &  86 &  70 & 12 \\
UFO      & 10{,}750{,}846 & 129\,MB   & 193 &  91 & 37 \\
Hasselt  & 19{,}615{,}568 & 235\,MB   & 385 & 368 & 24 \\
\bottomrule
\end{tabular}
\caption{Dataset properties extracted from PCD headers and bounding-box sampling.}
\label{tab:datasets}
\end{table}

\textbf{Corridor} is the smallest dataset. Its narrow Z extent (12\,m) relative
to its horizontal footprint is consistent with a ground-level indoor or
semi-outdoor corridor scan, where the sensor sweeps a bounded vertical range.

\textbf{UFO} is a mid-size outdoor scan ($\sim$10.8 million points). The
horizontal extents are moderate ($\sim$200\,m) and the Z span (37\,m) is the
largest of the three, reflecting a more varied vertical structure than the
other two datasets.

\textbf{Hasselt} is the largest dataset ($\sim$19.6 million points, 235\,MB).
Its footprint is roughly $385 \times 368$\,m but the Z extent is only 24\,m,
which is characteristic of a flat urban outdoor scan -- dense near facades and
the ground plane, sparse in open airspace. This dataset was used for all
stress-testing of memory usage and build time.

% ---------------------------------------------------------------------------
\section{Space Partitioning Algorithms}
\label{sec:algorithms}
% ---------------------------------------------------------------------------

All three implementations share a common base class (\texttt{BaseTree}) that
establishes a key memory design: throughout construction every algorithm works
exclusively with \textbf{integer indices} into the original
\texttt{Float32Array} rather than allocating one \texttt{Vector3} object per
point. Each index is a four-byte integer; coordinates are read inline as
\texttt{positions[idx*3]}, \texttt{positions[idx*3+1]},
\texttt{positions[idx*3+2]}. For Hasselt ($\sim$19.6 million points) the
construction index array consumes $\sim$78\,MB, compared to the $\sim$784\,MB
that a \texttt{Vector3[]} approach would require (roughly 40\,bytes per
JavaScript object versus 4 bytes per index). The global bounding box is
computed by the shared helper \texttt{\_computeBoundsFromPositions}, which
scans the \texttt{Float32Array} in a single pass and creates only two
\texttt{Vector3} objects (the AABB corners) regardless of point count. After
each tree is fully built, the algorithm-specific index structures and the
reference to the positions buffer are nulled, retaining only the node
structure in memory.

\subsection{Octree}
\textit{(Both authors)}

\subsubsection*{Design decisions}

The tree is built top-down: starting from a global bounding box, it is
recursively subdivided into 8 equal octants. A node stops splitting when
either \texttt{maxDepth} is reached \emph{or} its point count falls at or
below \texttt{maxPointsPerNode} -- whichever comes first.

The first design decision was to \textbf{force the root bounding box into a
perfect cube} by recentering it on the bounding-box centroid and expanding
all three half-extents to half of the longest side. This guarantees that
\emph{every voxel at a given depth has the same side length}, regardless of
the dataset's spatial extent. The consequences are threefold: the leaf-size
readout in the UI is meaningful because all leaves at the same depth are
equally sized; ray-AABB intersection tests benefit from uniform cell
dimensions that allow precomputed slab distances; and level-of-detail
rendering is straightforward -- querying a shallower depth uniformly
coarsens the whole scene at the same rate in every direction.

A second decision was \textbf{pre-allocating eight point-bucket arrays per
depth level} instead of creating them inside each recursive call. The na\"ive
approach would allocate eight new arrays per internal node; across thousands
of recursive calls on large datasets this generates significant
garbage-collection pressure. With the pre-allocated table each bucket is
reset to length zero via \texttt{.length = 0} and reused, keeping
allocation overhead constant regardless of tree size.

Each point is assigned to its octant using a \textbf{3-bit index}: bit~0
is set if the X coordinate $\geq$~\texttt{centerX}, bit~1 for Y, bit~2 for Z.
The same index drives the child bounding-box computation in $O(1)$ via bitwise
conditions, keeping the assignment and the geometry fully consistent without any
per-axis branching.

Only non-empty octants produce child nodes; technically, this makes the
implementation a \textbf{sparse octree} (up to 8 children per node, not
always 8). This is crucial for large, non-uniform point clouds such as
Hasselt, where most of the airspace above the city produces empty octants
-- skipping them avoids wasting memory and traversal time, and ensures only
occupied regions are rendered.

\subsubsection*{Complexity and memory}

\begin{itemize}
  \item \textbf{Build:} $O(n \log n)$ -- the split plane (geometric midpoint)
        is computed in $O(1)$; classifying each point into an octant requires
        three comparisons ($O(1)$ per point). Each of the $n$ points is touched
        once per depth level, across $O(\log n)$ levels.
  \item \textbf{Single-point update:} $O(\log n)$ -- follow the path from
        root to leaf and add/remove the point from the appropriate voxel;
        the split planes never change.
  \item \textbf{Memory:} $O(n)$ asymptotically -- at most $n$ leaf nodes at
        maximum depth; internal nodes null their index arrays once split.
        The dominant allocation during construction is the integer index array
        ($\sim$78\,MB for Hasselt), which is freed once the build completes,
        along with the positions reference and the pre-allocated bucket table.
        The retained node structure size scales with the number of occupied
        octants, which depends on the dataset's spatial distribution and the
        chosen depth.
\end{itemize}

\begin{figure}[H]
  \centering
  \includegraphics[width=0.48\linewidth]{figures/octree_depth_low.png}
  \hfill
  \includegraphics[width=0.48\linewidth]{figures/octree_depth_high.png}
  \caption{Octree on the Corridor dataset at a low depth (left) and a high
           depth (right). Color encodes the Z coordinate (blue = low,
           red = high). The wireframe outlines show voxel boundaries.}
  \label{fig:octree}
\end{figure}

% ---------------------------------------------------------------------------
\subsection{K-d Tree}
\textit{(Both authors)}
% ---------------------------------------------------------------------------

\subsubsection*{Design decisions}

The split plane is always placed at the \textbf{median} of the active
coordinate. Using the median rather than the geometric midpoint is the key
choice: it guarantees that each split divides the point set into two roughly
equal halves regardless of point density, keeping the tree balanced.

\paragraph{Axis-selection strategies.}
Three strategies are implemented and selectable at run time via the
\textbf{Split Mode} control in the viewer:

\begin{itemize}
  \item \textbf{Cycle} (default): the axis rotates X~$\to$~Y~$\to$~Z with
        depth (\texttt{depth~\%~3}). Simple and data-independent; all nodes
        at the same depth always split on the same axis.
  \item \textbf{Widest}: picks the axis whose bounding-box extent is largest
        at the current node, using the tight AABB already computed during
        construction. $O(1)$ per node; ensures the longest dimension is always
        halved first, producing more cube-like cells and reducing long thin
        partitions.
  \item \textbf{Variance}: picks the axis with the highest variance of the
        point coordinates -- the dominant diagonal entry of the covariance
        matrix, equivalent to what a full PCA would identify as the principal
        axis of spread when restricted to axis-aligned splits.
        Variance is estimated via \textbf{Welford's online algorithm} on a
        stride sample of up to 1\,024 points, keeping per-node cost $O(1)$ in
        practice while remaining numerically stable for coordinates of any
        scale. Unlike \textit{widest}, this strategy responds to the actual
        distribution of points rather than just the bounding geometry.
\end{itemize}

A module-level scratch \texttt{Vector3} is reused across all
\textit{widest} computations to avoid per-call heap allocations.

\IfFileExists{figures/kdtree_cycle.png}{%
\begin{figure}[H]
  \centering
  \includegraphics[width=0.32\linewidth]{figures/kdtree_cycle.png}
  \hfill
  \includegraphics[width=0.32\linewidth]{figures/kdtree_widest.png}
  \hfill
  \includegraphics[width=0.32\linewidth]{figures/kdtree_variance.png}
  \caption{K-d Tree split mode comparison on the Corridor dataset: Cycle (left),
           Widest (center), Variance (right). In Cycle mode all cells at a given
           depth share the same color; in Widest and Variance the colors are
           mixed, reflecting data-driven axis choices at each node.}
  \label{fig:kdtree_modes}
\end{figure}}{}

\paragraph{Median selection.}
Finding the median without a full sort is done with a
\textbf{3-way quickselect} (Dutch National Flag partitioning). This is
$O(n)$ average per node and handles duplicate coordinates correctly by grouping
equal-valued points into a pivot zone in a single pass, so repeated coordinates
in the cloud do not degrade performance to $O(n^2)$ as they would with a
na\"ive quicksort.

All points are stored in a \textbf{single shared array} that is partitioned
in-place during construction. Each node records a \texttt{[start, end)} range
into this array -- no per-node point copies are made. Bounding boxes are
computed from the actual point distribution after partitioning, giving tight
cuboids instead of the oversized AABB of the raw split range.

\subsubsection*{Balancing}

Because every split is at the median, the k-d tree is \textbf{structurally
balanced}: each partition divides the point set into two equal-count halves,
bounding the tree depth at $\lceil \log_2 n \rceil$. In practice, however,
the \texttt{maxPointsPerNode} leaf threshold means that a sparse region can
become a leaf several levels earlier than a dense region that keeps splitting.
The tree therefore remains \textbf{well-balanced} -- traversal depth stays
close to $O(\log n)$ -- but the strict guarantee that all leaves sit at the
same depth does not hold when the leaf threshold is active.

A degenerate edge case is caught explicitly: after quickselect, if the left
partition would be empty (\texttt{medianIndex == start}), the node is
finalized as a leaf immediately rather than recursing into a zero-size split.
This guards against configurations where the leaf threshold is zero and a
single remaining point would otherwise produce an empty left child.

A balanced tree is especially important for ray traversal: in a balanced tree
the worst-case traversal depth is $O(\log n)$, whereas an unbalanced tree can
degrade to $O(n)$.

\subsubsection*{Complexity and memory}

\begin{itemize}
  \item \textbf{Build:} $O(n \log n)$ -- quickselect is $O(n)$ per node;
        there are $O(\log n)$ levels; total work per level is $O(n)$.
  \item \textbf{Query / ray intersection:} $O(\log n)$ average due to balance.
  \item \textbf{Single-point update:} $O(\log n)$ for search; however,
        na\"ive insertion without rebalancing degrades the tree over time.
  \item \textbf{Memory:} $O(n)$ during construction (the shared
        \texttt{Int32Array} of $n$ indices, $\sim$78\,MB for Hasselt). After
        build both the index array and the positions reference are freed.
        The retained structure has at most $2^{\texttt{maxDepth}+1}-1$ nodes;
        for the depth-8 configuration used here that is $\leq$511 nodes
        ($\sim$100\,kB), so all three cached variants (Cycle, Widest, Variance)
        coexist in memory without retaining any copy of the point cloud.
\end{itemize}

\paragraph{Memory problem encountered.}
An initial prototype represented each point as a \texttt{Vector3} JavaScript
object. Each object carries a heap overhead of roughly 40--80\,bytes on top of
the 12 bytes of coordinate data; the construction array for Hasselt therefore
consumed $\sim$784\,MB--1.5\,GB, exceeding the browser heap limit and causing
an out-of-memory crash even for a single tree build.

The solution was to replace the \texttt{Vector3[]} array with a single
\texttt{Int32Array} of $n$ integer indices into the original \texttt{Float32Array}.
At 4\,bytes per index the construction array for Hasselt drops to $\sim$78\,MB.
After the tree is fully built both the index array and the positions reference
are nulled, leaving each cached variant at $\sim$100\,kB.
Additionally, the Variance axis-selection computation was capped at a stride
sample of 1\,024 points per node: without the stride, each node at the top of
the tree would scan all of its $O(n)$ points, making the per-node \emph{cost}
$O(n)$ rather than $O(1)$. The stride bounds it to at most 1\,024 iterations
per node regardless of subtree size. After these changes all three Hasselt
variants coexist in memory without issue.

\begin{figure}[H]
  \centering
  \includegraphics[width=0.48\linewidth]{figures/kdtree_corridor.png}
  \hfill
  \includegraphics[width=0.48\linewidth]{figures/kdtree_ufo.png}
  \caption{K-d Tree on the Corridor dataset (left) and the UFO dataset
           (right). The cuboids are non-uniform because the median split
           follows the actual point distribution, not a fixed midpoint.}
  \label{fig:kdtree}
\end{figure}

% ---------------------------------------------------------------------------
\subsection{BSP Tree}
\textit{(Antonio Romero Mart\'in)}
% ---------------------------------------------------------------------------

\subsubsection*{Design decisions}

A BSP tree splits space with \textbf{arbitrarily oriented planes}. The central
design challenge is choosing a good plane at each node. We evaluated three
strategies:

\begin{enumerate}
  \item \textbf{Random plane} -- fast but produces uneven splits.
  \item \textbf{Axis-aligned median} -- equivalent to k-d tree; does not
        exploit non-axis-aligned structure.
  \item \textbf{PCA-based plane} (Principal Component Analysis) -- the plane
        normal is the dominant eigenvector of the covariance matrix of the
        node's points; the plane is positioned at the \textbf{median
        projection} of the node's points onto that axis.
\end{enumerate}

We chose PCA because it adapts to the local geometry: in a corridor the
principal axis aligns along the corridor, producing long thin cells; in an
open area the plane follows the dominant spread of points. This reduces the
likelihood of degenerate splits compared to random or axis-aligned strategies.

The eigenvector is found by \textbf{power iteration} (up to 20 iterations on
the $3\times3$ covariance matrix, with early termination when the vector
magnitude falls below $10^{-10}$), which converges reliably for real-world
point clouds. The covariance matrix is computed in a \textbf{single pass}
using \textbf{Welford's online algorithm}: for each point the running mean is
updated first, then all six covariance entries are accumulated as
$\Delta_{\text{before}} \times (\text{value} - \mu_{\text{new}})$.
No scratch buffers are required; only six scalar accumulators are maintained
alongside the three mean values.

Once the principal axis $\mathbf{n} = (n_x, n_y, n_z)$ and centroid
$\mathbf{o}$ are known, every point is projected onto the axis:
$d_i = (\mathbf{p}_i - \mathbf{o}) \cdot \mathbf{n}$.
The \textbf{median projection} $d_{\text{med}}$ is found via a
\textbf{3-way quickselect} (Dutch National Flag) in $O(n)$ average, and the
split plane is placed at $\mathbf{o} + d_{\text{med}}\,\mathbf{n}$.
Points with $d_i \geq d_{\text{med}}$ go to the front child; the rest go to
the back child. This guarantees that each split divides the point set into two
equal-count halves, bounding the tree depth at $\lceil \log_2 n \rceil$ and
eliminating the $O(n^2)$ worst case of the centroid-based split.

A node stops splitting when either \texttt{maxDepth} is reached \emph{or} its
point count falls at or below \texttt{maxPointsPerNode}. If the PCA plane
produces a degenerate partition -- all points falling on the same side -- the
node is finalized as a leaf immediately to avoid infinite recursion on
numerically degenerate inputs.

\subsubsection*{Why BSP takes longer to build than an Octree}

An octree's split plane is always the three axis-aligned midplanes of the
current node -- computed in $O(1)$ regardless of point count. A BSP node
must \textbf{examine all its points} to compute the covariance matrix (Welford
pass), find the median projection (quickselect pass), and classify points
(one more pass): that is $O(n_\text{node})$ per node with a larger constant
than the Octree. With the median-projection split the tree is balanced and the
total work sums to $O(n \log n)$ -- the same asymptotic class as the Octree --
but the per-node constant is roughly $3\times$ higher. On the Hasselt dataset
($\sim$19.6 million points, 235\,MB) we measured BSP taking roughly $3\times$
longer than the Octree.

\subsubsection*{Complexity and memory}

\begin{itemize}
  \item \textbf{Build:} $O(n \log n)$ \emph{guaranteed} -- the median
        projection ensures each split divides the point set in half, bounding
        the tree depth at $\lceil \log_2 n \rceil$; each level processes $O(n)$
        points in total across three $O(n_\text{node})$ passes (Welford,
        quickselect, classify). The $O(n^2)$ worst case of the centroid-based
        split is eliminated.
  \item \textbf{Single-point update:} $O(\log n)$ to locate the correct leaf
        (balance is now guaranteed). Insertion is additionally expensive:
        adding a point can shift the local PCA plane, potentially invalidating
        the split for that node and requiring a rebuild of the affected subtree.
  \item \textbf{Memory:} $O(n)$ at all times -- during construction each
        recursive call allocates a local \texttt{Float32Array} of dot products
        for its node, freed as the call unwinds. Across the call stack at most
        ${\sim}2n$ elements exist simultaneously (${\sim}156\,$MB for Hasselt),
        less than the $3n$ (${\sim}234\,$MB) of the previous scratch-buffer
        approach. Unlike the k-d tree, the BSP tree \emph{retains} $O(n)$
        memory after build: each leaf stores a \texttt{Int32Array} of its point
        indices, needed by the visualizer for per-cell point recolouring. For
        Hasselt all leaf arrays together hold ${\sim}78\,$MB permanently,
        versus ${\sim}100\,$kB for the cached k-d tree.
\end{itemize}

\begin{figure}[H]
  \centering
  \includegraphics[width=0.48\linewidth]{figures/bsp_depth_low.png}
  \hfill
  \includegraphics[width=0.48\linewidth]{figures/bsp_depth_high.png}
  \caption{BSP Tree on the Corridor dataset at depth~2 (left) and depth~6
           (right). Each cell is assigned a distinct HSL hue; at low depth
           a few large regions are visible, while at high depth the
           partitioning follows the fine structure of the point cloud.}
  \label{fig:bsp_algo}
\end{figure}

% ---------------------------------------------------------------------------
\section{Viewer}
\label{sec:viewer}
% ---------------------------------------------------------------------------

\subsection{Rendering Architecture}
\textit{(Both authors)}

The viewer uses a \texttt{PerspectiveCamera} (75° FOV) with
\texttt{OrbitControls} for mouse-driven orbit. Pan and right-click are
disabled to keep navigation simple; rotation damping (factor~0.08) smooths
the motion without overshoot. On each dataset load the camera is automatically
repositioned via a bounding-box fit: a tight axis-aligned bounding box is
computed, the longest of its three dimensions is taken, and the camera distance
is set to $1.6\times$ that value so the entire point cloud fills the viewport.

Two services decouple concerns cleanly. \texttt{datasetService} owns point
cloud loading, height-ramp colorization, and camera fitting.
\texttt{partitionService} owns tree construction, per-variant caching, and
depth queries. Trees are built \textbf{on first use} and stored under a
composite key \texttt{algorithm::dataset::variant}. Switching algorithms,
changing depth, or toggling between the k-d tree's three split-mode variants
is instant once the corresponding tree is cached -- no rebuild is triggered on
subsequent interactions.

For Octree and k-d Tree, active nodes at the current depth are rendered as
\textbf{two \texttt{InstancedMesh} objects}: a semi-transparent solid box mesh
and a wireframe overlay, both backed by a single shared \texttt{BoxGeometry}.
Using \texttt{InstancedMesh} instead of one \texttt{Mesh} per node reduces
draw calls from thousands to one, which proved essential: without instancing,
rendering depth-8 Octree nodes ($\sim$8000 boxes) dropped the frame rate below
10~fps; with instancing it stayed above 60~fps.

\paragraph{Depth navigation.}
The \texttt{+} and \texttt{-} keys change the display depth. The tree is
queried for nodes exactly at depth $d$, and leaves shallower than $d$ are also
included so coarser resolutions still show all occupied regions. This traversal
is $O(\text{active nodes})$ per key press, while the expensive build runs only
once.

\subsection{Coloring}
\textit{(Both authors)}

\textbf{Octree} cells are colored by the Z coordinate of their bounding-box
center using a \textbf{height ramp}: HSL hue is linearly mapped from 0.7
(blue/purple, low~Z) to 0.0 (red, high~Z). This makes the vertical structure
of urban point clouds immediately visible and matches the conventional LiDAR
viewer palette.

\textbf{K-d Tree} cells are colored by the split axis that created them --
that is, the axis of the parent split that produced the boundary between a
node and its sibling. The three colors match the on-screen XYZ axis helper
exactly: \textcolor{red}{red} for X, \textcolor{green}{green} for Y, and
\textbf{blue} for Z. Because both children of a split receive the same axis
tag, sibling cells always share the same hue, making the direction of each
partition boundary immediately legible without any additional annotation.
With the \textit{Cycle} mode all cells at a given depth have the same color;
with \textit{Widest} and \textit{Variance} the colors at a given depth can
be mixed, directly reflecting the data-driven axis choices made at each node.

\paragraph{Wireframe.}
Each active node is rendered with a second \texttt{InstancedMesh} using
wireframe mode, making voxel boundaries visible without obscuring the interior
color. Solid fill and wireframe can be toggled independently via checkboxes in
the control panel.

\subsection{BSP Visualization}
\textit{(Antonio Romero Mart\'in)}

BSP cells are convex polytopes bounded by arbitrary planes, so axis-aligned
boxes are not a valid representation. We chose to \textbf{recolor the point
cloud in place}: each leaf cell is assigned a distinct hue by cycling evenly
across the HSL color wheel. This makes cell boundaries visible directly in the
data without needing to compute convex hulls (which would be expensive and
visually complex for thousands of cells).

Recoloring operates at the GPU buffer level. During construction each BSP leaf
records a compact \texttt{Int32Array} of the original geometry indices of its
points. When BSP mode is activated, the visualizer saves a full copy of the
Three.js \texttt{BufferGeometry} color attribute (the height-ramp colors
applied at load time) and then writes each leaf's HSL color directly into
the buffer at the stored indices. When the user switches away from BSP mode,
the saved copy is written back, restoring the baseline appearance without
reloading the dataset.

Additionally, the \textbf{split planes} are shown as semi-transparent quads
oriented to the PCA normal at each internal node up to the current depth
(Figure~\ref{fig:bsp}). Each quad is positioned at the median-projected origin of the split plane
(the point $\mathbf{o} + d_{\text{med}}\,\mathbf{n}$ stored in
\texttt{splitPlane}), scaled to the node's bounding-box size, and oriented via
quaternion rotation from the Z axis (the default normal of
\texttt{PlaneGeometry}) to the split plane normal. This gives an intuitive view of how space
is partitioned level by level.

\begin{figure}[H]
  \centering
  \includegraphics[width=0.48\linewidth]{figures/bsp_colored_cells.png}
  \hfill
  \includegraphics[width=0.48\linewidth]{figures/bsp_planes.png}
  \caption{BSP Tree: point cloud colored per leaf cell (left), and split
           planes shown as semi-transparent quads (right).}
  \label{fig:bsp}
\end{figure}

\subsection{UI and Dataset Loading}
\textit{(Siebe Ternest)}

\begin{figure}[H]
  \centering
  \includegraphics[width=0.36\linewidth]{figures/ui_panel.png}
  \caption{Control panel with algorithm selector, dataset buttons,
           depth display, leaf-size readout, and visibility toggles.}
  \label{fig:ui}
\end{figure}

The control panel (Figure~\ref{fig:ui}) provides algorithm and dataset
selection, solid/wireframe visibility toggles, and an XYZ axes toggle. Depth
and average leaf size are updated in real time. The \textbf{leaf size} readout
displays the mean cube root of active-node volumes,
$\bar{s} = \frac{1}{N}\sum_{i=1}^{N}\sqrt[3]{V_i}$, giving an intuitive
single-number summary of the current partitioning resolution that is
comparable across algorithms and datasets.

In addition to the three bundled datasets (UFO, Corridor, Hasselt), users can
drag and drop any \texttt{.pcd} file onto the viewer. The application runs
entirely in the browser; datasets are fetched over HTTP or read from the local
file system via the File API.

When the k-d Tree algorithm is selected, a \textbf{Split Mode} selector
appears in the control panel offering the three axis-selection strategies
(Cycle, Widest, Variance). Each variant is cached independently under a
composite key \texttt{algorithm::dataset::mode}, so switching between modes
after the first build is instant. The active mode is highlighted with the
same yellow accent used for dataset buttons, making the current strategy
immediately visible.

% ---------------------------------------------------------------------------
\section{Comparative Analysis}
\label{sec:analysis}
% ---------------------------------------------------------------------------

\subsection{Main Differences between the Three Trees}
\textit{(Both authors)}

\begin{center}
\begin{tabular}{@{}lccc@{}}
\toprule
 & \textbf{Octree} & \textbf{K-d Tree} & \textbf{BSP Tree} \\
\midrule
Split planes per node & 3 axis-aligned & 1 axis-aligned & 1 arbitrary \\
Children per node     & Up to 8        & Exactly 2      & Exactly 2 \\
Split position        & Fixed midpoint & Data median    & PCA median projection \\
Guaranteed balanced   & No             & Yes (in practice)      & Yes \\
Build cost (constant) & Low            & Medium         & High \\
Cell shape            & Cube           & Cuboid         & Convex polytope \\
\bottomrule
\end{tabular}
\end{center}

The Octree always cuts at the geometric midpoint along all three axes at once,
giving $O(1)$ split decisions. The k-d tree cuts one axis at a time at the
median, making it self-balancing at the cost of a linear scan. The BSP tree
uses an arbitrary plane adapted to the local point geometry, which produces the
most informative partitions but incurs the highest build cost.

\IfFileExists{figures/hasselt_octree.png}{%
\begin{figure}[H]
  \centering
  \includegraphics[width=0.48\linewidth]{figures/hasselt_octree.png}
  \hfill
  \includegraphics[width=0.48\linewidth]{figures/hasselt_kdtree.png}
  \caption{Hasselt dataset ($\sim$19.6 million points): Octree (left) and k-d Tree
           (right) at comparable depths. The height ramp and axis-color schemes
           make the structural differences visible at large scale.}
  \label{fig:hasselt}
\end{figure}}{}

\subsection{Benefits of an Octree over a BSP-Tree}
\textit{(Both authors)}

\begin{itemize}
  \item \textbf{Simpler, faster construction:} No covariance computation;
        the split planes are determined in $O(1)$.
  \item \textbf{Easy updates:} Inserting or removing a point only affects the
        nodes on its root-to-leaf path; split planes never change, so no
        structural rebuild is needed.
  \item \textbf{Predictable cell shapes:} All cells at a given depth are
        cubes of identical size, simplifying ray-AABB tests and frustum culling.
  \item \textbf{Implicit empty-space compression:} Only occupied octants
        create child nodes, so the sparse octree uses no memory for empty
        regions -- a significant advantage for large, non-uniform outdoor point
        clouds where most of the volume is empty air.
\end{itemize}

\subsection{Ray Intersection Queries}
\textit{(Both authors)}

All three trees support $O(\log n)$ average ray traversal by pruning subtrees
whose bounding volumes do not intersect the ray. In practice:

\begin{itemize}
  \item \textbf{Octree:} Ray-AABB intersection (slab method) is very fast
        with six comparisons per node; the regular grid at each level enables
        front-to-back ordering, allowing early termination once the first hit
        is found.
  \item \textbf{K-d Tree:} Each internal node test reduces to a single axis
        comparison (dot product with an axis vector), making it the cheapest
        per-node test. Balance guarantees $O(\log n)$ depth.
  \item \textbf{BSP Tree:} Each node requires a full dot product with an
        arbitrary normal, which is slightly more expensive per node than a
        single axis comparison. The median-projection split guarantees balance,
        so traversal depth is bounded at $O(\log n)$; however the higher
        per-node cost still makes it slower than the k-d tree in practice.
\end{itemize}

For ray casting, the \textbf{k-d tree} is generally the best choice: its
balance guarantee minimizes tree depth, and axis-aligned planes minimize
per-node cost.

\subsection{Dynamic Scenes}
\textit{(Both authors)}

For dynamic scenes with moving objects or streaming LiDAR, we would use an
\textbf{Octree}. Its regular structure makes insertion and deletion local
operations: a moving point only modifies the nodes on its root-to-leaf path,
and split planes never need updating. Unlike the k-d tree, which typically requires periodic full rebuilds to restore
balance after many insertions, the octree supports truly incremental updates
with no rebalancing step. The BSP tree's data-dependent planes make any update
expensive since the PCA result changes when the point distribution shifts,
potentially invalidating the split planes of entire subtrees.
\subsection{Urban Scenes}
\textit{(Both authors)}

For urban scenes we would use an \textbf{Octree}. While the k-d tree is also a
valid choice -- its median split adapts well to non-uniform point densities --
the octree has two practical advantages that matter specifically for urban
LiDAR data.

First, urban point clouds are extremely large (millions of points) and
spatially non-uniform: very dense near building facades and the ground plane,
almost empty in open air. A sparse octree compresses this empty space
implicitly and at multiple resolutions simultaneously, making it straightforward
to implement level-of-detail rendering by simply querying a shallower depth.
The k-d tree does not have a natural notion of uniform resolution levels,
since its cuboids vary in size and shape across the tree.

Second, urban mapping applications typically require repeated spatial queries
of the same type (e.g.\ nearest-neighbour search, frustum culling, occupancy
checks), all of which benefit from the octree's predictable, regular cell
structure. Practical LiDAR mapping tools such as OctoMap confirm this choice
as the industry standard for large-scale outdoor scenes.

% ---------------------------------------------------------------------------
\section{Conclusion}
\textit{(Both authors)}
% ---------------------------------------------------------------------------

We implemented three spatial partitioning structures -- Octree, k-d Tree, and
BSP Tree -- and integrated them into a single interactive browser-based viewer
for LiDAR point clouds, built with Three.js rather than raw OpenGL. This
allowed us to focus on algorithm design and rendering architecture without
low-level graphics boilerplate, at no cost to visual quality or interactivity.

The comparative analysis confirmed that no single tree dominates across all
criteria. The \textbf{octree} is the most practical general-purpose structure:
fast to build, easy to update incrementally, and naturally suited to
level-of-detail rendering of large urban scenes. The \textbf{k-d tree} is the
best choice for ray casting and nearest-neighbour queries, where its balance
guarantee and cheap per-node axis test minimize traversal cost. The
\textbf{BSP tree} produces the most geometry-adaptive partitions, but its
high PCA build cost and the fact that any point insertion invalidates the local
PCA planes make it less suitable for dynamic or real-time applications.

On the implementation side, three decisions had the largest practical impact.
\textbf{Replacing \texttt{Vector3} arrays with integer indices} during tree
construction reduced peak memory for Hasselt from $\sim$784\,MB to $\sim$78\,MB,
eliminating out-of-memory crashes and making all three k-d tree split-mode
variants cacheable simultaneously.
\textbf{InstancedMesh rendering} reduced thousands of draw calls to one,
keeping the frame rate above 60~fps at deep tree levels where per-node meshes
failed entirely.
\textbf{Tree caching} -- building each tree on first use and storing it under a
composite (algorithm, dataset, variant) key -- was the single most impactful UX
decision, making all interactions instant even on the largest tested dataset.
A fast-path binary PCD parser was also implemented: for files with the standard
\texttt{x~y~z~Float32} layout (all three bundled datasets), the position buffer
is created directly from the raw bytes after the header, eliminating $\sim$58.8
million \texttt{DataView} calls for Hasselt and reducing the post-download
parsing phase from seconds to milliseconds.

Overall, the project reinforced that the choice of spatial data structure is
inseparable from the intended query type, update frequency, and scene
characteristics -- a lesson that extends well beyond point cloud processing to
any geometry-heavy application.

\end{document}
